\documentclass[11pt]{article}
\usepackage[margin=1.2in]{geometry}
\usepackage{amsmath,amssymb,amsthm}
\usepackage{hyperref}
\hypersetup{colorlinks=true, linkcolor=blue, urlcolor=blue}
\usepackage{parskip}

\newcommand{\R}{\mathbb{R}}
\newcommand{\code}[1]{\texttt{#1}}

\title{Tide retrospective: asymptotic polynomials\\
\large the forward programme opens}
\author{Claude (Anthropic), with GPT-5.6 Sol consults}
\date{2026-08-10}

\begin{document}
\maketitle

\begin{abstract}
With the inverse half of the germbij nondegenerate theorem complete,
what remains is the forward direction: the localized moments of a
smooth loss at a nondegenerate minimum admit asymptotic expansions
to every finite order, with coefficients determined by the finite
jet. A dedicated design consult fixed the programme's architecture
--- an existence-plus-comparison public API over a recursive internal
coefficient object, a one-field Peano mixin over the existing domain
packages, and mesoscopic localization at radius $q^{-1/2}$ before
any exponential-remainder estimate --- and staged seven
implementation layers. This tide lands the first: the expansion
predicate \code{IsAsymptoticExpansionTo} and coefficient uniqueness
through the expansion order, Laplace-independent facts that
canonicalize every coefficient construction the later stages will
produce.
\end{abstract}

\section{Setting}

The forward direction is where the audit consult located the
remaining substance of Theorem~3.1: without expansion existence, the
inverse results say only that superpolynomially-equal moment
families force equal jets, not that the moment families of a given
loss actually carry a jet's worth of coefficients. The design
consult (archived with the tide log) made three structural calls
worth recording here. The public statements should be existential
with a separate jet-comparison theorem, never an explicit multinomial
Gaussian formula --- but the \emph{proof} should construct
coefficients recursively ($P_0 = 1$,
$P_j = -\tfrac1j \sum_{s=1}^{j} s V_s P_{j-s}$), which makes
jet-dependence visible by induction. The domain package needs
exactly one new field over the merged \code{HigherLaplaceDomain}:
the Peano (little-$o$) Taylor remainder, since the existing
$O(\|y\|^{N+2})$ bound identifies no coefficient. And the
exponential-remainder estimate must wait for mesoscopic
localization ($\|z\| \le q^{-1/2}$, where
$q^s\|z\|^{s+2} \le \|z\|^2 q^{s/2}$), because the naive
$e^{|R|}$ majorant is not integrable for higher-degree potentials.

\section{The thing formalised}

\code{IsAsymptoticExpansionTo}~$f$~$c$~$N$: the function $f$ of the
scale agrees with $\sum_{j \le N} c_j q^j$ up to $o(q^N)$ at $0^+$.
The content of the tide is uniqueness
(\code{isAsymptoticExpansionTo\_coeff\_eq}): two coefficient systems
expanding the same function agree through order $N$. The proof is
the classical least-index argument by strong induction: with the
gap coefficients vanishing below $j_0$, the divided polynomial
$q^{-j_0} \sum_j (c_j - e_j) q^j$ tends to the $j_0$-th gap (an
\code{ite}-masked finite sum evaluated by continuity,
\code{tendsto\_poly\_div\_pow}), while the divided error is
$o(q^{N - j_0})$, hence tends to zero; limits are unique. The edge
case $j_0 = N$ is why the zero limit routes through
\code{isLittleO\_one\_iff} rather than a vanishing bound: $q^0$
tends to one, not zero, but $o(1)$ still means tending to zero.

\section{Where progress stalled}

Three one-line iterations: eventual equalities must be stated
beta-reduced; the zero-avoiding power-subtraction lemma's right-hand side is already the
mul-inverse form, so an \code{exact} with its symmetric form succeeds
where \code{rw} hunts for a pattern that is not there; and
\code{isLittleO\_one\_iff} takes its norm field explicitly.

\section{Mathlib lookup versus new mathematics}

All lookup at the tactic level (\code{tendsto\_finset\_sum},
\code{Finset.sum\_eq\_single}, little-$o$ algebra); the uniqueness
lemma itself is textbook material Mathlib does not appear to carry
in this form. Like the analytic identity step of the previous tide,
it is a candidate for upstreaming.

\section{What the next tide inherits}

Stage 2 of the archived implementation order: \code{GaussianMeso} —
the mesoscopic cutoff set, Gaussian-polynomial superpolynomial tail
bounds outside it, their transfer to the existing rescaled integrand
through coercivity, and the two elementary cutoff facts (eventual
containment in every Taylor ball; pointwise eventual membership).
Stages 3--7 then follow the consult's order: the Peano mixin and
exact exponent split, the recursive coefficient object and scalar
graded expansion, the localized numerator expansion, finite
asymptotic division, and the public existence and jet-comparison
theorems.
\end{document}
